% GENERATED FILE. Edit canonical lessons, then run npm run generate:books.
% canonical-source-hash: fnv1a64:ccf9b4297974b080
% canonical-lessons: TA-C401-natattunar, TA-C401-niruttam, TA-C401-kulunku, TA-C401-tallatu, TA-C401-cavari

\chapter{On a Crowded Bus}
\label{ch:ta-on-a-crowded-bus}
\hlchaptermodality{401}

\begin{chapteropening}
\textbf{By the end of this chapter:} \emph{I can talk about a bus conductor, a stop, a jolting ride and staggering.}

Complete the last lesson of chapter 401: I can talk about a bus conductor, a stop, a jolting ride and staggering.
\end{chapteropening}

\section[\texorpdfstring{\ta{நடத்துநர்} (naṭattunar)}{நடத்துநர் (naṭattunar)}]{\ta{நடத்துநர்} (naṭattunar) — a bus conductor}

\label{lesson:TA-C401-natattunar}



\begin{quote}
Before the new one: say the Tamil for to postpone, then the Tamil for to switch.
\end{quote}

\subsection*{You'll want to know: \ta{நடத்துநர்}}
\textbf{\ta{நடத்துநர்}} — \emph{naṭattunar} — \textquotedblleft{}a bus conductor\textquotedblright{}.

\begin{quote}
\ta{நடத்துநர்} \ta{எனக்குச்} \ta{சீட்டு} \ta{கொடுத்தார்}.
\end{quote}

\begin{quote}
\emph{naṭattunar eṉakkuc cīṭṭu koṭuttār.}
\end{quote}

\begin{quote}
The conductor gave me a ticket.
\end{quote}

The ending -\ta{நர்} makes a doer from a verb: \ta{நடத்து}, to run or conduct, gives \ta{நடத்துநர்}.

\begin{cousinweb}
\textbf{\ta{நடத்துநர்}} (\emph{naṭattunar}) is built from \ta{நடத்து}, \textquotedblleft{}to conduct, to make go\textquotedblright{} (DEDR 3582, the causative of \ta{நட}, to walk), and the ending \emph{-nar}, \textquotedblleft{}one who\textquotedblright{}.
\end{cousinweb}

\begin{grammarlens}[title={What you've built}]
One more word for getting around.
\end{grammarlens}

\subsection*{Your turn}
\begin{itemize}
\raggedright
  \item \emph{Say it:} \emph{naṭattunar}
  \item \emph{Say it:} \emph{naṭattunar}, once more
  \item \emph{Say it:} \emph{naṭattunar eṉakkuc cīṭṭu koṭuttār.} — The conductor gave me a ticket.
  \item \emph{Recall:} say the Tamil for to postpone, then the Tamil for to switch, then say \emph{naṭattunar} again
\end{itemize}

\begin{tcolorbox}[breakable,colback=teal!4,colframe=teal!35!black,title={Before you move on}]
What does \ta{நடத்துநர்} mean? (\textquotedblleft{}A bus conductor\textquotedblright{}.) Say it once more.
\end{tcolorbox}
\section[\texorpdfstring{\ta{நிறுத்தம்} (niṟuttam)}{நிறுத்தம் (niṟuttam)}]{\ta{நிறுத்தம்} (niṟuttam) — a bus stop; a stop}

\label{lesson:TA-C401-niruttam}



\begin{quote}
Before the new one: say the Tamil for to switch, then the Tamil for a bus conductor.
\end{quote}

\subsection*{You'll want to know: \ta{நிறுத்தம்}}
\textbf{\ta{நிறுத்தம்}} — \emph{niṟuttam} — \textquotedblleft{}a bus stop; a stop\textquotedblright{}.

\begin{quote}
\ta{அடுத்த} \ta{நிறுத்தத்தில்} \ta{நான்} \ta{இறங்குவேன்}.
\end{quote}

\begin{quote}
\emph{aṭutta niṟuttattil nāṉ iṟaṅkuvēṉ.}
\end{quote}

\begin{quote}
I will get off at the next stop.
\end{quote}

\ta{பேருந்து} \ta{நிறுத்தம்} is a bus stop; before a case it becomes \ta{நிறுத்தத்தில்}.

\begin{cousinweb}
\textbf{\ta{நிறுத்தம்}} (\emph{niṟuttam}) is built from \ta{நிறுத்து}, \textquotedblleft{}to stop\textquotedblright{} (DEDR 3675, the family of \ta{நில்}, to stand), with the noun ending \emph{-am}.
\end{cousinweb}

\begin{grammarlens}[title={What you've built}]
One more word for getting around.
\end{grammarlens}

\subsection*{Your turn}
\begin{itemize}
\raggedright
  \item \emph{Say it:} \emph{niṟuttam}
  \item \emph{Say it:} \emph{niṟuttam}, once more
  \item \emph{Say it:} \emph{aṭutta niṟuttattil nāṉ iṟaṅkuvēṉ.} — I will get off at the next stop.
  \item \emph{Recall:} say the Tamil for to switch, then the Tamil for a bus conductor, then say \emph{niṟuttam} again
\end{itemize}

\begin{tcolorbox}[breakable,colback=teal!4,colframe=teal!35!black,title={Before you move on}]
What does \ta{நிறுத்தம்} mean? (\textquotedblleft{}A bus stop; a stop\textquotedblright{}.) Say it once more.
\end{tcolorbox}
\section[\texorpdfstring{\ta{குலுங்கு} (kuluṅku)}{குலுங்கு (kuluṅku)}]{\ta{குலுங்கு} (kuluṅku) — to jolt; to shake}

\label{lesson:TA-C401-kulunku}



\begin{quote}
Before the new one: say the Tamil for a bus conductor, then the Tamil for a bus stop.
\end{quote}

\subsection*{You'll want to know: \ta{குலுங்கு}}
\textbf{\ta{குலுங்கு}} — \emph{kuluṅku} — \textquotedblleft{}to jolt; to shake\textquotedblright{}.

\begin{quote}
\ta{மோசமான} \ta{சாலையில்} \ta{பேருந்து} \ta{குலுங்கியது}.
\end{quote}

\begin{quote}
\emph{mōcamāṉa cālaiyil pēruntu kuluṅkiyatu.}
\end{quote}

\begin{quote}
The bus jolted on the bad road.
\end{quote}

Past (it …): \ta{குலுங்கியது}, \emph{kuluṅkiyatu}.

\ta{குலுங்கு} is to shake by itself; \ta{குலுக்கு} is to shake something, as in \ta{கை} \ta{குலுக்கு}.

\begin{cousinweb}
\textbf{\ta{குலுங்கு}} (\emph{kuluṅku}) is native Dravidian (DEDR 1806); the same article holds \ta{குலுக்கு}, \textquotedblleft{}to shake something\textquotedblright{}.
\end{cousinweb}

\begin{grammarlens}[title={What you've built}]
One more word for getting around.
\end{grammarlens}

\subsection*{Your turn}
\begin{itemize}
\raggedright
  \item \emph{Say it:} \emph{kuluṅku}
  \item \emph{Say it:} \emph{kuluṅku}, once more
  \item \emph{Say it:} \emph{mōcamāṉa cālaiyil pēruntu kuluṅkiyatu.} — The bus jolted on the bad road.
  \item \emph{Recall:} say the Tamil for a bus conductor, then the Tamil for a bus stop, then say \emph{kuluṅku} again
\end{itemize}

\begin{tcolorbox}[breakable,colback=teal!4,colframe=teal!35!black,title={Before you move on}]
What does \ta{குலுங்கு} mean? (\textquotedblleft{}To jolt; to shake\textquotedblright{}.) Say it once more.
\end{tcolorbox}
\section[\texorpdfstring{\ta{தள்ளாடு} (taḷḷāṭu)}{தள்ளாடு (taḷḷāṭu)}]{\ta{தள்ளாடு} (taḷḷāṭu) — to stagger; to totter}

\label{lesson:TA-C401-tallatu}



\begin{quote}
Before the new one: say the Tamil for a bus stop, then the Tamil for to jolt.
\end{quote}

\subsection*{You'll want to know: \ta{தள்ளாடு}}
\textbf{\ta{தள்ளாடு}} — \emph{taḷḷāṭu} — \textquotedblleft{}to stagger; to totter\textquotedblright{}.

\begin{quote}
\ta{பேருந்து} \ta{திடீரென்று} \ta{நின்றதால்} \ta{நான்} \ta{தள்ளாடினேன்}.
\end{quote}

\begin{quote}
\emph{pēruntu tiṭīreṉṟu niṉṟatāl nāṉ taḷḷāṭiṉēṉ.}
\end{quote}

\begin{quote}
I staggered because the bus stopped suddenly.
\end{quote}

Past (I …): \ta{தள்ளாடினேன்}, \emph{taḷḷāṭiṉēṉ}.

The second part is \ta{ஆடு}, to move or sway, as in \ta{நடனமாடு}.

\begin{cousinweb}
\textbf{\ta{தள்ளாடு}} (\emph{taḷḷāṭu}) is native Dravidian: DEDR 3128 lists it, \textquotedblleft{}to totter\textquotedblright{}, built with \ta{ஆடு}, \textquotedblleft{}to move, to sway\textquotedblright{}.
\end{cousinweb}

\begin{grammarlens}[title={What you've built}]
One more word for getting around.
\end{grammarlens}

\subsection*{Your turn}
\begin{itemize}
\raggedright
  \item \emph{Say it:} \emph{taḷḷāṭu}
  \item \emph{Say it:} \emph{taḷḷāṭu}, once more
  \item \emph{Say it:} \emph{pēruntu tiṭīreṉṟu niṉṟatāl nāṉ taḷḷāṭiṉēṉ.} — I staggered because the bus stopped suddenly.
  \item \emph{Recall:} say the Tamil for a bus stop, then the Tamil for to jolt, then say \emph{taḷḷāṭu} again
\end{itemize}

\begin{tcolorbox}[breakable,colback=teal!4,colframe=teal!35!black,title={Before you move on}]
What does \ta{தள்ளாடு} mean? (\textquotedblleft{}To stagger; to totter\textquotedblright{}.) Say it once more.
\end{tcolorbox}
\section[\texorpdfstring{\ta{சவாரி} (cavāri)}{சவாரி (cavāri)}]{\ta{சவாரி} (cavāri) — a ride}

\label{lesson:TA-C401-cavari}



\begin{quote}
Before the new one: say the Tamil for to jolt, then the Tamil for to stagger.
\end{quote}

\subsection*{You'll want to know: \ta{சவாரி}}
\textbf{\ta{சவாரி}} — \emph{cavāri} — \textquotedblleft{}a ride\textquotedblright{}.

\begin{quote}
\ta{குழந்தைகள்} \ta{யானைச்} \ta{சவாரி} \ta{செய்தார்கள்}.
\end{quote}

\begin{quote}
\emph{kuḻantaikaḷ yāṉaic cavāri ceytārkaḷ.}
\end{quote}

\begin{quote}
The children went for an elephant ride.
\end{quote}

\ta{சவாரி} \ta{செய்} is to go for a ride; \ta{யானைச்} \ta{சவாரி}, an elephant ride.

\begin{cousinweb}
\textbf{\ta{சவாரி}} (\emph{cavāri}) is a loan from Persian \emph{savārī}, \textquotedblleft{}riding, a mount\textquotedblright{}, which reached Tamil through Hindustani.
\end{cousinweb}

\begin{grammarlens}[title={What you've built}]
One more word for getting around.
\end{grammarlens}

\subsection*{Your turn}
\begin{itemize}
\raggedright
  \item \emph{Say it:} \emph{cavāri}
  \item \emph{Say it:} \emph{cavāri}, once more
  \item \emph{Say it:} \emph{kuḻantaikaḷ yāṉaic cavāri ceytārkaḷ.} — The children went for an elephant ride.
  \item \emph{Recall:} say the Tamil for to jolt, then the Tamil for to stagger, then say \emph{cavāri} again
\end{itemize}

\begin{tcolorbox}[breakable,colback=teal!4,colframe=teal!35!black,title={Before you move on}]
What does \ta{சவாரி} mean? (\textquotedblleft{}A ride\textquotedblright{}.) Say it once more.
\end{tcolorbox}
